\scanpage{177}\mspage{32}%
% archive sheet number: 177
% C original page filename: H134-05_134-2_scan177.tex
où ``pris à la lettre'' $p''$ est maintenant une inclusion, mais
``moralement'' c'est \inserted{tjrs} une fibration à fibres très
\unsure{polies|jolies} [de dim finie notamment, voire des variétés
non singulières -- c'est donc beaucoup mieux que ce que nous donne
le yoga de Cartan-Serre ``toute application continue est équivalente
à une fibration''\ldots] C'est ce dernier diagramme qui est susceptible
d'une construction directe, purement algébrique, dans le cadre des
schémas, lorsque l'on a développé la construction du voisinage tubulaire étale.

Le point où je veux en venir, c'est que la considération
\inserted{de la cat} des faisceaux (ou \inserted{complexes d'iceux, ou}
$n$-champs \unsure{..|\ldots}) $\unsure{\mathcal P|P}$-constructibles
sur un $X$, où $\mathcal P$ est une stratification ``équisingulière''
donnée, revient dans le cas particulier actuel à la connaissance d'un
diagramme de \uline{types d'homotopie} ordinaires
% The small added parenthesis is physically above the continuation word ordinaires;
% no attachment caret fixes a more exact position within this qualification.
\inserted{(ou de pro-types, si on revient à la top. étale)}%
% Isolated near-vertical mark at approximately (1332,1565), above/right of the
% closing parenthesis; reproduced without assigning a prime or punctuation role.
\raisebox{1.1ex}[0pt][0pt]{\parbox{0pt}{\figure{\begin{tikzpicture}[overlay]
\draw[line width=0.2pt] (0,0) -- (0.005,-0.06);
\end{tikzpicture}}}}
\diagram{\begin{tikzcd}[ampersand replacement=\&,column sep=large,row sep=large]
\unsure{\Delta|\triangle} \arrow[r,"\unsure{J|j}"] \arrow[d,"p"'] \& X_1\\
X_0 \&
\end{tikzcd}}
% Isolated short dark stroke at approximately (316,1851), to the left of and
% below the diagram, above the first post-diagram prose line. Origin uncertain.
\noindent\raisebox{0pt}[0pt][0pt]{\parbox{0pt}{\figure{\begin{tikzpicture}[overlay]
\draw[line width=0.3pt] (0.85,0.55) .. controls (0.87,0.52) and (0.85,0.51) .. (0.87,0.48);
\end{tikzpicture}}}}%
\unsure{en|et} prenant des systèmes de coeff locaux (ou de $n$-champs
loc. constants) sur ces sommets
\inserted{$X_0,X_1$ \deleted{$\unsure{\Delta|X}$}}
reliés par des hom de compatibilité du type
$p^*(F_0)\to J^*\unsure{(}F_1)$. Ce serait un exercice amusant
(que je n'\unsure{ai} pas encore fait) de
